\section{MVCC and deletion for sequence fragments}
\label{sec:pg-storage-versions}

I discussed MVCC support when talking about Postgres Snapshots in chapter about BioseqDB custom scans~\ref{sec:scan-row-stability}.
In custom scan case support is rather simple and provided by the Postgres engine itself. It only required me to correctly use it.
Storing fragments is entirely separate problem to visibility in columnar storage, because multiple queries can operate on locators, while some other process can 
modify sequence data in the same time. 

Deleting sequence require doing two operations:
\begin{itemize}
  \item Deleting row for fragment, which under the hood will mark the row as deleted in the row-mask metadata. 
  \item As fragment is separate abstraction layer than columnar storage, I also store versions of rows in locators, so deletion just needs to invalidate that version (epoch)
\end{itemize}

Each epoch has stored separate row in the catalog table (so it's one entry per sequence version):
\noindent\begin{minipage}{\linewidth}
\begin{lstlisting}[style=bioseq-storage-sql,
  basicstyle=\footnotesize\ttfamily,
  breaklines=true,breakatwhitespace=true,
  caption={Heap catalog for one immutable sequence version. Invalidation updates \texttt{state}.},
  label={lst:storage-source-version}]
CREATE TABLE bioseq_internal.sequence_source_version (
  source_epoch bigint PRIMARY KEY, -- immutable version id, must be > 0
  user_relid oid,                  -- parent user table
  database_oid oid,                -- parent database
  tablespace_oid oid,              -- parent tablespace
  main_relfilenode oid,            -- parent file node
  main_storage_id bigint,          -- parent TAM storage generation
  column_attno smallint,           -- 1-based sequence column
  locator_version smallint,        -- locator format, currently 3
  bucket_idx smallint,             -- length bucket 0..34
  shard_idx integer,               -- shard inside the bucket
  fragment_relid oid,              -- fragment table, NULL if empty
  fragment_storage_id bigint,      -- fragment TAM storage generation
  first_sub_row_no bigint,         -- row of fragment 0
  num_fragments bigint,            -- consecutive fragment rows
  total_bp_size bigint,            -- sequence length in symbols
  alphabet_code smallint,          -- packed alphabet
  storage_unit_byte_width smallint,-- unit width in bytes
  padding_letters integer,         -- unused letters in last packed unit
  source_content_checksum bigint,  -- CRC-64 of decoded symbols
  state smallint                   -- 0 unpublished, 1 ready, 2 invalid
);
\end{lstlisting}
\end{minipage}

The specific epoch's \texttt{state} follows very simple linear transitions:
\begin{itemize}
  \item \textbf{Unpublished} ($0$). An insert has allocated an epoch (sequence version) and written fragments.
   Locator is then filled with state $0$ which is default one. It means that the locator in the user table is not yet written
   (so the fragment is there, but there's no saved locator that points to it)
  \item \textbf{Ready} ($1$). Locator was written to a table, so I update the epoch state to mark that this sequence version is used somewhere
   (not doing that might lead to dangling fragments that occupy storage but have no existing locator). This state is crucial to prevent a situation in which locator write fails or is rolled back, but fragment was written.
  \item \textbf{Invalid} ($2$). Locator cannot read this row version anymore, the underlying fragment shouldn't be accessed.
    This happens when user either calls \texttt{DELETE} or \texttt{UPDATE}.
\end{itemize}

Note that the catalog contains far more information that is neccessary. I decided to save CRC hashes to make sure the fragment is not malformed during operations
or shard coordinates shared by the locator itself. As each row version requires separate entry this limits the maximum number of sequences stored in the BioseqDB, however
entries are cheap so I didn't incorporate it into estimate calculations~\ref{eq:storage-human-genome-table-cap}.

\begin{figure}[htbp]
\centering
\begin{tikzpicture}[x=1cm,y=1cm,font=\footnotesize]
  \draw[storage arrow,black!55] (1.75,4.55)--(12.2,4.55) node[right] {time};
  \draw[dashed,black!55] (7.25,.1)--(7.25,4.9);
  \node[above,align=center] at (7.25,4.9) {writer commits\\the update};
  \node[anchor=east] at (1.5,3.7) {reader A};
  \draw[storageBlue,line width=2pt] (2.0,3.7)--(10.8,3.7);
  \fill[storageBlue] (2.0,3.7) circle (2pt);
  \node[above,text=storageBlue] at (4.3,3.7) {snapshot before update};
  \node[below,text=storageBlue] at (9.25,3.7) {later read: ACGT};
  \node[anchor=east] at (1.5,2.55) {reader B};
  \draw[storageOrange,line width=2pt] (8.1,2.55)--(11.7,2.55);
  \fill[storageOrange] (8.1,2.55) circle (2pt);
  \node[above,text=storageOrange] at (9.9,2.55) {snapshot after update};
  \node[below,text=storageOrange] at (9.9,2.55) {read: TGCA};
  \node[storage cell,fill=storageBlue!12,minimum width=4.65cm,
    minimum height=1.15cm] (old) at (4.3,.85)
    {SQL id 1; internal row 42\\epoch 7: ACGT};
  \node[storage cell,fill=storageOrange!12,minimum width=4.65cm,
    minimum height=1.15cm] (new) at (10.0,.85)
    {SQL id 1; internal row 57\\epoch 8: TGCA};
  \node[anchor=west] at (1.85,-.2)
    {After commit, both sets of fragment bytes remain stored.};
\end{tikzpicture}
\caption{Postgres snaphot mechanism allows to check visivility of rows in MVCC model, but it's not sufficient on its own.
Let's assume that reader A and reader B are reading the same sequence before and after other process commited update.
Using fragments epochs versioning system means that read A will see older locator via snapshot, while reader B will
get newer locator row with changed epoch. Both of them can complete reads as both epochs are still valid. 
}
\label{fig:storage-update-snapshots}
\end{figure}

Listing~\ref{lst:storage-repeatable-read} gives SQL code for this example:

\begin{lstlisting}[style=bioseq-storage-sql,float=htbp,
  label={lst:storage-repeatable-read},
  caption={A repeatable-read transaction keeps the old sequence version.
    Initially, row id 1 begins with ACGTACGT.}]
-- Session A
BEGIN ISOLATION LEVEL REPEATABLE READ;
SELECT bioseq.to_text(seq[0:8]) FROM reads WHERE id = 1;
-- ACGTACGT

-- Session B, while A remains open
BEGIN;
UPDATE reads
SET seq = repeat('TGCA', 4096)::bioseq.seq('dna4', 32)
WHERE id = 1;
COMMIT;

-- Session A
SELECT bioseq.to_text(seq[0:8]) FROM reads WHERE id = 1;
-- ACGTACGT
COMMIT;
SELECT bioseq.to_text(seq[0:8]) FROM reads WHERE id = 1;
-- TGCATGCA
\end{lstlisting}
